\begin{bookbox}{inset}{Engineering at the frontier of authority}
\label{inset-frontier-authority}
The need to make expertise count is clearest when engineers advise someone else with authority to proceed. Engineers may not make the final decision, but they remain responsible for making the engineering case legible: what is understood, what can be assured, what remains uncertain, and what risk a decision would accept.

Engineering history contains painful demonstrations of this distinction. Before Challenger, Morton Thiokol engineers recommended against launching at the temperatures expected on January 28, 1986. Management reversed that recommendation, but the engineering objection survived in the record, allowing the Rogers Commission to reconstruct both the technical disagreement and the decision process that followed it\citep{rogers1986}. Columbia illustrates a different failure: engineers raised concerns about the foam strike and sought additional imagery, but important assumptions and uncertainties did not reach the Mission Management Team with sufficient force or clarity. The Columbia Accident Investigation Board subsequently called for an independent technical authority and safety-assurance organization with authority over safety decisions\citep{caib2003}.

Volkswagen\textquotesingle s diesel-emissions scandal brings the problem into software. Engineers implemented software that recognized emissions tests and changed vehicle behavior; the subsequent legal record made it possible to reconstruct what engineers and managers knew, what they concealed, and who participated in the decisions. Volkswagen engineer James Liang ultimately pleaded guilty to conspiracy and was sentenced to 40 months in federal prison\citep{doj-vw-liang}. The disclosures of Frances Haugen and, later, former Facebook engineering director Arturo Béjar bring the pattern entirely into software systems. Facebook\textquotesingle s own researchers studied harms experienced by young users, while Béjar later testified about internal measurements of such harms and his efforts to communicate them to senior executives\citep{senate-haugen2021, senate-bejar2023}. Whatever judgment one makes about the appropriate product or policy response, these records made a more basic question answerable: what did the organization know about the behavior of its system, who knew it, and what did those with authority do with that knowledge?

These cases suggest a simple professional test. When consequential harm occurs, the engineering record should show what engineers knew, what they did not know, what they warned about, and who chose to proceed. This is not merely defensive documentation. Making knowledge, uncertainty, disagreement, and evidence explicit is part of how engineering informs consequential decisions and makes those decisions answerable.

The distinction matters increasingly as software agents become capable of consequential action. Retaining engineering control does not require engineers to make every decision an agent makes. It requires them to establish the relevant bounds: what authority has been delegated, what properties are assured within it, what evidence supports those claims, what departures can be detected, and where the limits of that assurance lie. The engineering challenge is to extend the frontier over which meaningful authority can be exercised as the capabilities and autonomy of the system expand.

Sometimes the available evidence will not establish the desired bound. Engineers are not always the final decision-makers. Executives, military commanders, regulators, political leaders, and societies may have legitimate reasons to proceed despite engineering uncertainty: the consequences of inaction may themselves be severe, and competitive or geopolitical pressures may make accepting substantial risk rational. Engineering\textquotesingle s responsibility is to preserve the distinction. A decision to accept an inadequately bounded risk does not make the risk bounded. The record should make clear what engineering could establish, what it could not, and who possessed the authority to proceed anyway.

That answerability matters after the decision as well. Society needs identifiable people and institutions from whom an account can be demanded and to whom consequences can attach. Professional licenses can be revoked, certifications withdrawn, organizations sanctioned, and individuals held legally responsible. A machine cannot become the terminus of that chain of responsibility. Delegating a decision to a machine does not delegate the responsibility that justified the human or institutional authority to make it.

Better engineering environments can expand the domain over which engineers responsibly delegate action. At this frontier, their value is measured by what people and institutions can understand, constrain, assure, and govern---and by whether they can answer for the consequences.
\end{bookbox}
